\chapter{Introduction}
\label{chap:introduction}

Imagine a world where ordinary glasses can discreetly record conversations and capture images or videos without your knowledge or consent. With the rise of augmented reality (AR) glasses, this is the reality we are entering. These devices are part of the development of the Internet of Things (IoT), alongside technologies such as voice assistants, smartwatches, and smart home gadgets. While they offer many conveniences for our daily lives, they also raise privacy concerns. Incidents of baby monitors being turned into spy devices or doorbells sharing video with third parties showcase the potential risks of these technologies. Unlike standard IoT devices, AR glasses introduce unique privacy concerns due to their ability to continuously record surroundings and transmit visual data and audio information from the wearer's perspective to remote servers. For instance, a recent report of Meta Ray-Ban glasses being used for facial recognition without consent \citet{IXRAY2025} illustrates how such devices could violate privacy and increase surveillance risks.

While existing research has explored privacy risks, focusing on a range of IoT devices, there remains a scarcity of studies examining the privacy implications of specific wearable AR devices. This thesis addresses this gap by analyzing the Meta Ray-Ban glasses, an AR wearable IoT device developed by Meta Platforms, Inc. (formerly Facebook) from Menlo Park, CA, USA. 

Studying the Ray-Ban Meta glasses presented several challenges. The proprietary nature of the glasses, in which manufacturers limit access to internal components, software, and data processing mechanisms, makes it difficult to understand and examine their data collection and processing mechanisms. This lack of transparency made it challenging to evaluate privacy risks effectively. Furthermore, the encryption protocols used in the network traffic, such as Transport Layer Security (TLS) and security techniques such as certificate pinning, protect the content of transmitted data. TLS encrypts the transmitted content, ensuring that unauthorized parties cannot read it, while certificate pinning prevents interception of adversaries by rejecting connections to unverified servers. Although these measures enhance security, they also make it challenging to analyze the specific information being sent to external servers or third parties. 

In light of these challenges, this thesis applies a combination of traffic and app analysis to examine the privacy implications of the Ray-Ban Meta glasses. By analyzing network traffic captured under a range of usage scenarios and investigating its companion app, Meta View, this thesis aims to identify and understand what data is transmitted or possibly collected, its destinations, and its potential effects on user and bystander privacy. This focused approach contributes to a more complete understanding of the unique privacy risks of wearable AR devices and provides a basis for developing effective mitigation strategies.  

The remainder of this thesis is organized as follows. Chapter \ref{chap:Background} presents the background and motivation for studying AR glasses, including an overview of relevant privacy and security concepts. Chapter \ref{chap:relatedWork} reviews related work in the field, focusing on prior works on AR devices and their privacy concerns. Chapter \ref{chap:deviceOverview} outlines the hardware and software components of the Ray-Ban Meta glasses. Chapter \ref{chap:methodologyExperiment} describes the methodology for data capture and analysis as well as their limitations. Chapter \ref{chap:appAnalysis} provides a static analysis of the companion app, Meta View, exploring permission requests and potential data flows. Chapter \ref{chap:networkAnalysis} discusses the network-level analysis for traffic captures from the experiments, while chapter \ref{chap:bluetoothAnalysis} analyzed the Bluetooth communication. Chapter \ref{chap:decryption} details the decryption attempts made and the analysis of the endpoints successfully decrypted. Chapter \ref{chap:discussion} discusses the key findings of the privacy risks and challenges encountered. Finally, chapter \ref{chap:conclusionFutureWork} presents the conclusion and proposes future work in this field.
